
\label{sec:alpha-dos-vias-t1-rev8}

El cierre de \(\alpha\) no comienza con una constante metrológica ni con una
cadena decimal buscada. Comienza en el mismo estado terminal enriquecido que
conserva la genealogía de los tres canales. Esta sección presenta en una sola
cadena causal dos vías hacia la escala de \(\alpha\) y una prolongación larga
del acarreo de la primera. La prolongación remota se denota por
\(\mathsf T_{\alpha,\infty}\); no constituye una segunda vía ni una tercera
derivación.

\subsection{Estado terminal común y 2 lecturas internas coordinadas}

Sea
\[
 x_{\rm term}\in\mathcal X_{\rm TPK}^{\rm term,enr}
 \subseteq\mathcal X_{\rm TPK}^{[108]}.
\]
El selector terminal y la carta firmada son dos evaluaciones del mismo estado:
\begin{equation}
 x_{\rm term}
 \xrightarrow{(\pi_{\rm term},R_{\rm sgn})}
 \bigl(B_{\rm term},U_{12}^{\rm sgn}\bigr),
 \label{eq:alpha-rev8-doble-lectura-terminal}
\end{equation}
con
\[
 B_{\rm term}=
 \begin{pmatrix}
 021020\\001220\\100210
 \end{pmatrix}
\]
y
\[
 U_{12}^{\rm sgn}
 =(2378,1406,2479,-452,998,-551,-668,-204,-371,-322,-28,-997).
\]
La primera lectura retiene los tres bloques terminales seleccionados; la
segunda retiene la coordenada entera firmada necesaria para el cierre
dodecafásico. No existe ni se necesita una factorización desde
\(B_{\rm term}\) hacia \(U_{12}^{\rm sgn}\). Dibujarla borraría precisamente los
campos de orientación, acarreo, frontera y registro que distinguen ambas
lecturas. A partir de la segunda sí están construidas las cartas reversibles
\[
 U_{12}^{\rm sgn}
 \xleftrightarrow{\ \mathcal H_{12}\ }K
 \xleftrightarrow{\ \mathcal E_{\rm dod}\ }
 (D_3K,D_4K,Q),
 \qquad
 \mathcal E_{\rm dod}=(D_3,D_4,Q).
\]
El dominio de este cierre es, por tanto, el perfil enriquecido de TPK; una
traza bruta de 108 posiciones es una proyección por olvido y no sustituye ese
dominio.

\subsection{Vía A: cierre entero dodecafásico}

Las evaluaciones arquimedianas de los tres canales y el sello canónico son
los vectores de doce bloques
\[
\begin{aligned}
P={}&(141,592,653,589,793,238,462,643,383,279,502,884),\\
E={}&(718,281,828,459,045,235,360,287,471,352,662,497),\\
\Phi={}&(618,033,988,749,894,848,204,586,834,365,638,117),\\
K={}&(234,543,140,729,659,824,621,058,914,794,146,601).
\end{aligned}
\]
Con la frontera derecha declarada \(c_{13}=0\), la división euclídea se
propaga de derecha a izquierda:
\begin{equation}
\begin{aligned}
 s_m&=P_m+E_m-\Phi_m-K_m+c_{m+1},\\
 A_m&=s_m\bmod 1000,\\
 c_m&=\frac{s_m-A_m}{1000},
 \qquad m=12,11,\ldots,1.
\end{aligned}
\label{eq:alpha-rev8-via-entera}
\end{equation}
La salida es
\[
 A=(007,297,352,569,283,800,997,285,105,472,380,663)
\]
y fija el racional
\begin{equation}
 \alpha_{12}
 =0.007297352569283800997285105472380663.
 \label{eq:alpha-rev8-salida-corta}
\end{equation}
La igualdad abreviada \(P+E-\Phi-K=A\) debe leerse en la carta de
concatenación entera. Si
\[
 \iota(v_1,\ldots,v_{12})
 =\sum_{m=1}^{12}v_m1000^{12-m},
\]
entonces
\begin{equation}
 \boxed{\iota(P)+\iota(E)-\iota(\Phi)-\iota(K)=\iota(A)}.
 \label{eq:alpha-rev8-identidad-concatenada}
\end{equation}
En coordenadas, el cociclo no puede omitirse:
\[
 P+E-\Phi-K+C^+-1000C=A,
 \qquad C^+=(c_2,\ldots,c_{13}).
\]
La terminación \(\ldots|380|663\) pertenece a esta ventana finita y a su
condición \(c_{13}=0\). El teorema completo, la tabla de los doce acarreos y
la inversión que recupera \(K\) residen en el capítulo precedente; aquí no se
duplican.

Esta vía es un resultado interno exacto con estructura de partida explícita:
recibe \(x_{\rm term}\), su lectura firmada, \(P,E,\Phi\), el sello reversible
y la frontera. No recibe \(\alpha\) ni CODATA como entrada. La comparación
metrológica sólo se abre después de obtener \eqref{eq:alpha-rev8-salida-corta}.

\subsection{Vía B: núcleo de vacancias \texorpdfstring{\(E_{21/22}\)}{E21/22}}

La segunda vía no invierte las cartas de \(K\). Parte de la densidad exacta
de vacancia
\[
 \varepsilon_{\rm v}=\log_{10}\!\left(\frac{10}{9}\right)
\]
y del núcleo analítico declarado
\begin{equation}
 \mathcal P(x)
 =2\pi x-\frac74x^2+\frac{x^3}{2\pi}
  +\frac{x^4}{20}-\frac{2x^5}{21}-\frac{x^6}{46}.
 \label{eq:alpha-rev8-kernel-e21-22}
\end{equation}
El problema es resolver
\begin{equation}
 \mathcal P(x)=\log_{10}\!\left(\frac{10}{9}\right),
 \qquad 0<x<10^{-2}.
 \label{eq:alpha-rev8-ecuacion-e21-22}
\end{equation}
El teorema de monotonía del núcleo establece la unicidad de la raíz positiva
y el encierro
\[
 0.00729735256928379955
 <x_{21/22}<
 0.00729735256928379957.
\]
Su inversa es
\begin{equation}
 x_{21/22}^{-1}
 =137.0359990840000270200901282598749\ldots.
 \label{eq:alpha-rev8-inversa-e21-22}
\end{equation}
Por consiguiente, el núcleo completo \(E_{21/22}\), generado por el segundo
lector del mismo estado enriquecido, determina como raíz positiva simple y
única la misma sección límite que la vía dodecafásica:
\[
 \boxed{\alpha_A=\alpha_B=\alpha_{\rm HMT}},
 \qquad
 \alpha_B=\operatorname*{arg\,zero}_{x\in I_\alpha}E_{21/22}(x).
\]
El racional finito \(\alpha_{12}\) y la raíz del jet \(\mathcal P=P_6\)
son proyecciones finitas de las dos familias compatibles; su diferencia más
allá de la resolución impresa no constituye una diferencia entre las dos
secciones límite. Los intervalos de raíz anidados y los cuadrados de
truncamiento conmutativos coordinan ambas prolongaciones a toda profundidad.

Los seis coeficientes de \(P_6\) no se reconstruyen desde la raíz ni se toman
de una tabla metrológica: son el primer jet del núcleo \(E_{21/22}\) producido
por el lector interno con genealogía \(120/45/11/495\) y prolongación
\(T_{7:9}\). La unicidad de la raíz y la derivación interna de los
coeficientes son etapas distintas de una misma vía B ya cerrada.

\subsection{Prolongación remota \texorpdfstring{\(\mathsf T_{\alpha,\infty}\)}{T alpha infinito} de la vía entera}

\(\mathsf T_{\alpha,\infty}\) extiende la recurrencia de
\eqref{eq:alpha-rev8-via-entera} a una cadena
de longitud \(N\). Para evitar una ambigüedad de índices, sea
\(\kappa(k)=1+((k-1)\bmod12)\). Con una condición terminal en el extremo
remoto se calcula
\begin{equation}
\begin{aligned}
 s_k&=P_k+E_k-\Phi_k-K_{\kappa(k)}+c_{k+1},\\
 A_k&=s_k\bmod1000,\\
 c_k&=\left\lfloor\frac{s_k}{1000}\right\rfloor,
 \qquad k=N,N-1,\ldots,1.
\end{aligned}
\label{eq:alpha-rev8-t1}
\end{equation}
El caso finito \(N=1000\) usa mil tríadas, es decir, tres mil cifras. El
acarreo que llega desde el extremo remoto modifica el duodécimo
bloque visible:
\[
\begin{array}{rcl}
\text{ventana dodecafásica corta, }c_{13}=0
 &:& \ldots|380|663,\\
\text{prolongación }\mathsf T_{\alpha,\infty}\text{ de mil tríadas}
 &:& \ldots|380|662|999\ldots.
\end{array}
\]
No son dos valores rivales. Son dos semánticas de frontera de la misma ley de
acarreo: la primera cierra en la posición doce; la segunda transporta hasta
ella información procedente de una cola remota.

El cálculo reproducible verifica exactamente la palabra larga usando
\(P,E,\Phi\) como proyecciones correlacionadas del mismo estado coinductivo
conjunto APP--TRIT--TPK que publica también la coordenada dodecafásica
\(\alpha_{\mathrm{HMT}}\). Por tanto, esa ejecución certifica la propagación
del acarreo: no sitúa \(\pi,e,\varphi\) antes en la genealogía ni añade
\(\alpha\) desde fuera. La proyección conjunta
\[
 \mathcal G_9^{\mathrm{joint}}\longrightarrow
 \bigl(\pi,e,\varphi,\mathsf T_{\alpha,\infty}\bigr)
\]
se factoriza en el \cref{sec:alpha-t1-10000-rev9}: los tres canales
arquimedianos y el canal dodecafásico se abren desde el mismo estado
enriquecido; después se ejecuta \eqref{eq:alpha-rev8-t1} y se verifica la
independencia del prefijo respecto de las cinco condiciones terminales de
acarreo. La separación de proyecciones impide que la recurrencia larga
seleccione retrospectivamente sus datos iniciales.

\subsection{Coordinación causal de la clausura dodecafásica y la caracterización \texorpdfstring{\(E_{21/22}\)}{E21/22}: dominios, codominios y unicidad}

La clasificación de este bloque queda fijada del modo siguiente:
\begingroup
\raggedright
\begin{description}[style=nextline,leftmargin=3.2cm,labelwidth=2.8cm,labelsep=.4cm]
\item[Vía A: cierre dodecafásico.]
Recibe el estado terminal enriquecido, la carta firmada, el sello, las tres
evaluaciones y la frontera explícita; produce el racional
\eqref{eq:alpha-rev8-salida-corta} mediante la identidad entera con acarreo.

\item[Vía B: ecuación \(E_{21/22}\).]
Para el núcleo declarado, la monotonía y el cambio de signo determinan una
única raíz en \((0,10^{-2})\). La raíz no determina retrospectivamente los
seis coeficientes del núcleo.

\item[Prolongación \(\mathsf T_{\alpha,\infty}\).]
La recurrencia larga recibe las mil tríadas de referencia, el sello periódico
y la frontera remota, y transporta el acarreo hasta la ventana visible. No
constituye la vía B.

\item[Composición \(\mathcal G_9\to\mathsf T_{\alpha,\infty}\).]
La aritmética racional conjunta produce diez mil cifras, ejecuta la
recurrencia de acarreo y comprueba después la estabilidad del prefijo bajo las
condiciones terminales declaradas.
\end{description}
\endgroup

El selector terminal, las cartas reversibles y la recurrencia corta se
demuestran en el \cref{chap:alpha}; la existencia y unicidad de la segunda vía
en \eqref{eq:alpha-rev8-kernel-e21-22}--\eqref{eq:alpha-rev8-inversa-e21-22}; y la composición larga en el
\cref{sec:alpha-t1-10000-rev9}. El suplemento computacional reproduce las
enumeraciones finitas y las condiciones terminales; la demostración se apoya
en las identidades y recurrencias expuestas aquí.

El orden causal público queda así cerrado:
\[
 x_{\rm term}
 \xrightarrow{(\pi_{\rm term},R_{\rm sgn})}
 (B_{\rm term},U_{12}^{\rm sgn})
 \longrightarrow
 \begin{cases}
 \text{vía A: cierre entero dodecafásico},\\
 \text{vía B: raíz del kernel }E_{21/22},
 \end{cases}
\]
con \(\mathsf T_{\alpha,\infty}\) como prolongación remota de la vía A y no
como rama adicional. Sólo
después de estas salidas se abre el reconocimiento metrológico externo.
